\chapter{Problem Statement and Motivation}

\section{Urban Traffic Bottlenecks and Intersection Failure Modes}
Traffic congestion in metropolitan centers represents a major impediment to economic productivity, environmental sustainability, and public safety. In rapidly expanding tech corridors such as Bengaluru, India, vehicular density consistently exceeds road capacity, resulting in notorious bottleneck nodes. The \textbf{Central Silk Board Junction}---a vital nexus connecting Electronic City, HSR Layout, BTM Layout, and Koramangala---frequently experiences clearance queues extending several kilometers.

Conventional municipal traffic management exhibits three fundamental failure modes:
\begin{enumerate}[leftmargin=*]
    \item \textbf{Rigid Fixed-Time Cycles}: Traditional electro-mechanical or simple micro-controller signal systems operate on predetermined phase timers (e.g., 45 seconds per approach regardless of instantaneous arrival rates). When one approach faces heavy commuter influx while another remains deserted, green time is squandered on empty asphalt, precipitating gridlock.
    \item \textbf{Manual Police Override Strain}: Traffic enforcement personnel must physically intervene, relying on subjective visual estimation to override signals manually. This introduces fatigue, inconsistent phase balancing, and delayed responses to approaching emergency responders.
    \item \textbf{Commuter Blindness and Upstream Information Asymmetry}: Commuters approaching the junction lack granular, corridor-specific signal phase status and congestion transparency, preventing proactive route selection and smooth braking.
\end{enumerate}

\section{The Role of Computer Vision and Adaptive Optimization}
Edge-deployed computer vision offers a non-intrusive, cost-effective alternative to costly subterranean inductive loop sensors or pneumatic road tubes. By processing optical feeds from standard high-definition traffic monitoring cameras, deep learning models can:
\begin{itemize}[leftmargin=*]
    \item Enumerate vehicles discretely across multiple classes (passenger cars, motorcycles, heavy freight trucks, and mass-transit buses).
    \item Compute instantaneous corridor queue densities in real time without physical road disruption.
    \item Identify anomalous vehicular behavior such as signal violations and counter-flow navigation.
\end{itemize}

\section{System Boundaries and Project Purpose}
BengaluruTrafficAI was engineered as an end-to-end full-stack software prototype to demonstrate how modern edge perception, adaptive optimization algorithms, and responsive reactive web architectures intersect.

\begin{conceptbox}[Scope Separation: Research Prototype vs. Municipal Actuation]
It is critical to distinguish the system's demonstrated software architecture from physical municipal deployment. The project does not directly actuate high-voltage physical traffic light hardware on Bengaluru streets, nor does it access restricted municipal CCTV feeds. Instead, it processes standardized generic urban sample video in its local perception mode and broadcasts simulated fluctuations in its cloud-hosted demonstration mode. The primary architectural objective is to establish a verified, test-covered, and decoupled software platform capable of ingesting streaming detections and executing robust signal optimization algorithms.
\end{conceptbox}
